\subsection{Inverse form of a bilinear or sesquilinear form.}

Let A be a ring, E a left A-module, F a right A-module, and \(\Phi\) a bilinear form on \(\mathbf{E} \times \mathbf{F}\) . It is assumed here that the associated maps of \(\Phi\) , which will be denoted \(s\) and \(d\) , are bijective. Then the product map \((s, d)\) is a bijection from \(\mathbf{E} \times \mathbf{F}\) onto \(\mathbf{F}^* \times \mathbf{E}^*\) , and defines, by transport of structure, a bilinear form \(\hat{\Phi}\) on \(\mathbf{F}^* \times \mathbf{E}^*\) . This therefore satisfies

\[
\begin{array}{r l} \hat {\Phi} (y ^ {\prime}, x ^ {\prime}) = & \Phi (s ^ {- 1} (y ^ {\prime}), d ^ {- 1} (x _ {i} ^ {\prime})) \\ = & \left\langle s ^ {- 1} (y ^ {\prime}), x ^ {\prime} \right\rangle = \left\langle d ^ {- 1} (x ^ {\prime}), y ^ {\prime} \right\rangle \qquad (x ^ {\prime} \in \mathrm{E} ^ {*}, y ^ {\prime} \in \mathrm{F} ^ {*}). \end{array}\tag{27}
\]

DEFINITION 8. --- Let \(\Phi\) be a bilinear form on E × F whose associated maps s and d are bijective. The bilinear form \(\hat{\Phi}\) on F* × E* defined by (27) is called the inverse form of \(\Phi\) .

Now let \(\hat{s}\) and \(\hat{d}\) be the linear maps from F* to E** and from E* to F** associated on the left and on the right to \(\hat{\Phi}\) . Since, for \(x' \in \mathrm{E}^*\) and \(y' \in \mathrm{F}^*\) , we have by definition

\[
\hat {\Phi} (y ^ {\prime}, x ^ {\prime}) = \left\langle y ^ {\prime}, \hat {d} (x ^ {\prime}) \right\rangle = \left\langle x ^ {\prime}, \hat {s} (y ^ {\prime}) \right\rangle
\]

Comparing with (27), one sees that the linear form \(\hat{d}(x')\) on F* is equal to that defined by the element \(d^{-1}(x')\) of F. It follows that the composite map \(\hat{d} \circ d\) is the canonical map from F to its bidual F**, and that this map is bijective since \(d\) and \(\hat{d}\) (the latter by transport of structure) are bijective; hence, if one canonically identifies F with F**, one has \(\hat{d} = d^{-1}\) . Similarly, E is canonically identified with E**, and the canonical map from E to E** is \(\hat{s} \circ s\) and one has \(\hat{s} = s^{-1}\) . It follows from this that the inverse form of \(\hat{\Phi}\) is \(\Phi\) .

Now consider a ring A equipped with an antiautomorphism J, two left A-modules E and F, and a right sesquilinear form \(\Phi\) for J on E \(\times\) F such that the associated maps of \(\Phi\) , which will be denoted \(s\) and \(d\) , are bijective. Define a map \(\hat{\Phi}\) of \(\mathbf{F}^{*} \times \mathbf{E}^{*}\) into A by the first equation (27). This map satisfies, by (24) (no. 6), the relation

\[
\hat {\Phi} (y ^ {\prime}, x ^ {\prime}) = \left\langle s ^ {- 1} (y ^ {\prime}), x ^ {\prime} \right\rangle = \left\langle d ^ {- 1} (x ^ {\prime}), y ^ {\prime} \right\rangle^ {J} \quad (x ^ {\prime} \in E ^ {*}, y ^ {\prime} \in F ^ {*}).\tag{28}
\]

The map \(\hat{\Phi}\) is evidently \(\mathbf{Z}\) -bilinear; moreover, we have, for \(a\) , \(b\) in A, \(x' \in \mathrm{E}^*\) and \(y' \in \mathrm{F}^*\) and by virtue of the definitions of \(s\) and \(d\) ,

\[
\widehat {\Phi} (y ^ {\prime} a, x ^ {\prime} b) = \Phi (a ^ {j} s ^ {- 1} (y ^ {\prime}), b ^ {j ^ {\prime}} d ^ {- 1} (x ^ {\prime})) = a ^ {j} \widehat {\Phi} (y ^ {\prime}, x ^ {\prime}) b;
\]

therefore \(\hat{\Phi}\) is a left sesquilinear form for J (no. 2) on \(\mathbf{F}^{*} \times \mathbf{E}^{*}\) .

DEFINITION 9. --- Let \(\Phi\) be a right-sesquilinear form for J on E × F, whose associated maps s and d are bijective. The left-sesquilinear form \(\hat{\Phi}\) for J on F* × E* is called the inverse form of \(\Phi\) .

We leave it to the reader to define and study the inverse form of a left sesquilinear form. This inverse form is a right sesquilinear form.

Let \(\hat{s}\) and \(\hat{d}\) be the associated maps of \(\hat{\Phi}\) ; by (26) (no. 6) we have

\[
\hat {\Phi} (y ^ {\prime}, x ^ {\prime}) = \left\langle y ^ {\prime}, \hat {d} (x ^ {\prime}) \right\rangle^ {\mathbf {J}} = \left\langle x ^ {\prime}, \hat {s} (y ^ {\prime}) \right\rangle .\tag{29}
\]

Because of the fact that \(s\) is bijective, and from the equality \(\langle s^{-1}(y'), x'\rangle = \langle y', \hat{d}(x') \rangle^{\mathfrak{J}}\) which results from (28) and (29), we deduce that \(\hat{d}\) is bijective; therefore \(\hat{d} \circ d\) is bijective. But the equality \(\langle d^{-1}(x'), y' \rangle = \langle y', \hat{d}(x') \rangle\) which results from (28) and (29), shows that \(\hat{d} \circ d\) is the canonical map from F into its bidual F**. One sees similarly that \(\hat{s}\) is bijective and that \(\hat{s} \circ s\) is the canonical map from E into E**. Therefore, if one identifies E** with E and F** with F by means of these canonical maps, we have \(\hat{s} = s^{-1}\) , \(\hat{d} = d^{-1}\) , and \(\Phi\) is the inverse form of \(\hat{\Phi}\) .

With the same notations and hypotheses, let a be an invertible element of the center of A. Then the maps associated with the form \(a\Phi\) are, by (23) (resp. (24)), equal to \(a.d\) and \(a.s\) (resp. \(a^{j'}\) .s), hence are bijective. It thus follows from (27) that the inverse form of \(a\Phi\) is \(a^{-1}\widehat{\Phi}\) (resp. \((a^{j'})^{-1}\widehat{\Phi}\) ).

